\chapwithtoc{Introduction}\label{chap:introduction}


Code optimization is a fundamental aspect of software development in high-performance computing (HPC). Modern HPC systems have complex architectures with hierarchies of memory and processing units, making a careful choice of data layouts in memory and traversal orders in algorithms crucial for achieving efficient execution.

A difference between a na\"{\i}ve implementation and an optimized one can lead to orders of magnitude improvements in execution time, runtime resources utilization, and energy efficiency~\cite{hennessy2019computer,drepper2007every}. Thus, if an algorithm is not properly tuned to the specific system it is running on, its execution can unnecessarily waste a tremendous amount of time and energy.

% TODO: There are many algorithm characteristics that affect the performance; we focus on data layouts and traversal orders in both sequential and parallel algorithms.
There are many factors that affect the performance of an algorithm on modern HPC architectures, such as algorithmic complexity, choice of data types and instructions, or arithmetic intensity of computations. We focus on memory access patterns, which are fundamental to optimizing the performance of algorithms, as many algorithms are bottlenecked by memory bandwidth and latencies~\xxx{(CITE)}. This issue is further pronounced with increasing size of processed data, as well as with the scale of parallelism in modern HPC systems, which requires data to be communicated and synchronized between different processing units, leading to additional overheads.

We explore approaches to transforming data layouts and traversal orders that are orthogonal to the other aspects of algorithm design, and thus can be applied to a wide variety of algorithms without compromising further optimization opportunities.

% TODO: memory has performance limitations
Optimizing data access patterns is important mainly because of the performance limitations of memory units in modern HPC architectures, which have specialized hardware features that improve the performance of data access for specific patterns of access. Most often, such patterns involve accessing data in a contiguous and sequential way and, if the algorithm makes repeated accesses to the same data, those accesses should happen within a short time. Similarly, certain access patterns allow for better exploitation of parallelism. Thus, careful choice of data layouts and their corresponding traversal orders is crucial for achieving efficient execution of algorithms on modern HPC architectures.

Even simple data structures such as matrices can be stored in a variety of layouts in memory, such as row-major, where elements of each row are stored contiguously; or column-major, where elements of each column are stored contiguously. If an algorithm traverses a given matrix row by row, the row-major layout will tend to access the data in almost sequential order, while the column-major layout will tend to access the data in a more scattered way.

Many features of modern HPC architectures benefit from accessing data in a contiguous way --- for example, CPU caches store multiple consecutive data elements together in a single addressable unit called a cache line for faster access during computation, making the row-major layout significantly more efficient for row-wise traversal, compared to the column-major layout. Modern systems also provide many independent computation units that can operate in parallel --- for these, the data layouts and traversal orders should be designed in such a way that each unit can access the data it needs with minimal interference with other units, minimizing data dependencies between data accessed by different units.

Since contemporary CPU and GPU architectures have many features that benefit from specific data access patterns, and algorithms typically access multiple data structures in multiple ways, the choice of appropriate data layouts and traversal orders is challenging and often requires empirical experimentation to find the best combination of them. Furthermore, different architectures may require different data access patterns for optimal performance, making it difficult to design an algorithm in such a way that it can efficiently run on different architectures without needing to be significantly redesigned for each one.

Many tools and frameworks have been developed to help design and optimize applications for different hardware architectures, such as polyhedral model-based optimizers, autotuning frameworks, and domain-specific languages for various areas of scientific computing~\cite{bondhugula2008pluto,ansel2014opentuner,ragan2013halide,baghdadi2019tiramisu}. However, such tools are often limited in their applicability to algorithms and do not expose optimization transformations to users in clear and intuitive ways.

Our contributions explore abstraction design to separate the concerns of algorithmic logic and data layout and traversal optimizations, and the potential of Large Language Models (LLMs) to assist with algorithm design. Both groups of our contributions are motivated by the goal of making algorithm optimization more accessible to users and thus improving the performance of scientific applications on modern HPC architectures.

We have introduced \emph{layout-agnostic} and \emph{traversal-agnostic} abstractions that allow users to express algorithmic logic independently of specific data layouts and traversal orders, enabling users to design and optimize them separately.
Our approach is based on composition of small, reusable, and semantically intuitive building blocks that can be used to express a wide variety of data layout and traversal transformations. We have shown its applicability to general scientific computing on distributed CPU and GPU architectures.

We have also studied the potential of LLMs to assist with code optimization, focusing on their ability to design efficient code for GPUs, and the interoperability of LLMs with the abstractions for user-directed code optimization. For both of these aspects, we have identified key challenges and opportunities for future research in this area.

\paragraph{Thesis structure.} The rest of this thesis is structured as follows. \Cref{chap:background} outlines how data layouts and traversal orders affect the performance of HPC code on modern CPU and GPU architectures. \Cref{chap:optimization-methods} presents an overview of existing methods for code optimization in this domain. \Cref{chap:abstractions} presents our work on abstractions for user-directed code optimization and \cref{chap:llms} presents our work on LLM-assisted code optimization. The thesis concludes with \cref{chap:conclusion}, which summarizes our contributions \xxx{and outlines future research opportunities in this area}.
